% Generated draft: requires mathematical and source review.
\documentclass[11pt]{article}
\usepackage[margin=25mm]{geometry}
\usepackage{fontspec}
\setmainfont{Latin Modern Roman}
\usepackage{amsmath,amssymb,mathtools,unicode-math}
\setmathfont{Latin Modern Math}
\usepackage{enumitem,needspace,xurl,hyperref}
\setlength{\parindent}{0pt}
\setlength{\parskip}{5pt}
\setlength{\emergencystretch}{3em}
\newcommand{\N}{\mathbb{N}}
\newcommand{\Z}{\mathbb{Z}}
\newcommand{\Q}{\mathbb{Q}}
\newcommand{\R}{\mathbb{R}}
\newcommand{\C}{\mathbb{C}}
\newcommand{\F}{\mathbb{F}}
\newcommand{\A}{\mathbb{A}}
\newcommand{\PP}{\mathbb{P}}
\newcommand{\E}{\mathbb{E}}
\newcommand{\eps}{\varepsilon}
\newcommand{\dd}{\,\mathrm{d}}
\begin{document}
\section*{UnsolvedMath ID: 30004706}
\textbf{Source problem code:} catalogue entry
\subsection{Definitions and mathematical statement}
For a K3 surface \(S\), the Hilbert scheme \(S^{[n]}=\operatorname{Hilb}^{n}(S)\) parametrizes zero-dimensional subschemes of \(S\) of length \(n\), where \(n\) is a nonnegative integer. This notation does not specify which values of \(n\) the source intends or how \(n\) relates to its other parameters.

The supplied record denotes an enumerative series associated with such Hilbert schemes by
\[
F_{g,m}(\alpha;\gamma_1,\ldots,\gamma_N).
\]
It does not define this series, the meanings or allowed ranges of \(g,m,N\), or the spaces containing \(\alpha\) and the insertions \(\gamma_i\).

The requested task is to establish the following three assertions identified collectively as Conjectures 2--4 in the cited source, with the hypotheses and quantifiers intended there:
\begin{enumerate}
\item The series \(F_{g,m}(\alpha;\gamma_1,\ldots,\gamma_N)\) has the conjectured quasi-Jacobi-form structure, with the index and weight prescribed by the source.
\item These series satisfy the conjectured multiple-cover formula expressed using formal Hecke operators.
\item These series satisfy the conjectured holomorphic-anomaly equation.
\end{enumerate}
The supplied record contains neither the predicted index and weight nor the multiple-cover formula, the formal Hecke operators, or the holomorphic-anomaly equation. It also omits the enumerative theory, generating variables, normalization conventions, underlying field, and geometric hypotheses needed to define the series. Consequently, the three assertions above specify the intended targets but cannot be expanded into fully quantified mathematical statements from the supplied evidence alone. The precise formulations of the cited Conjectures 2--4 must be supplied before this draft can constitute a complete conjecture statement.
\subsection{Short English statement}
The problem asks for three structural properties of enumerative generating series for Hilbert schemes of points on a K3 surface: quasi-Jacobi-form structure with prescribed index and weight, a multiple-cover formula using formal Hecke operators, and a holomorphic-anomaly equation. The supplied record omits the definitions, hypotheses, and formulas needed to state these assertions precisely.
\subsection{Sources}
\begin{itemize}
\item Local UnsolvedMath registry, supplied record with permanent ID \(30004706\), title ``Quasi-Jacobi Structure of K3 Hilbert-Scheme Invariants,'' exact source problem code \texttt{catalogue entry}; collection: Oberwolfach Reports Open Problems. No individual proposer is supplied; the proposed year is recorded as 2021.
\item Supplied report citation: \emph{Enumerative Geometry of Surfaces} (2022), Oberwolfach Reports, pp.~1463--1518, DOI \texttt{10.4171/owr/2021/28}. \url{https://doi.org/10.4171/owr/2021/28}. The record identifies Conjectures 2--4 as the source of the three requested assertions but does not reproduce them. The citation gives 2022, while the DOI contains 2021 and the proposed year is 2021; these supplied details are retained without resolution.
\item Supplied literature source labeled ``Oberdieck'': \url{https://people.kth.se/~georgo/note/HilbHAE.pdf}. No document title, version, theorem numbers, or more specific attribution is supplied.
\item Literature assessment dated exactly \texttt{2026-08-22}: the supplied curation labels the problem ``partially solved'' and states that the quasi-Jacobi, multiple-cover, and holomorphic-anomaly assertions have proofs in genus zero and in low-marking cases, while the stated general conjectures remain open. This is supplied dated curation, not an independently verified assessment. The background explicitly requires independent verification of time-sensitive status and cited claims.
\end{itemize}
\end{document}
